% Appendix A: Lemma Catalogue.  Every listing is copied verbatim from
% tamarin/final/lightning.spthy or tamarin/final/lightning_aux.spthy (proof scripts omitted).
\section{Lemma Catalogue}\label{app:catalogue}

This appendix renders every lemma of Table~\ref{tab:groups} in natural language and gives its
exact \tamarin{} statement, so that the table can be checked against the theories. We omit two
kinds of statements: that of \code{Skipped_Requires_Compromise}, which is printed as
Lemma~\ref{lem:skip}, and those of the seven helper lemmas of the general theorem, which we list
by name. The last subsection covers the lemmas of the auxiliary models.

Unless stated otherwise, every lemma below is verified in \code{lightning.spthy}, the main model
(channel layer, multi-hop routing in which one generic rule \code{Forward_HTLC} serves every hop
of a route of any length, and the wormhole rule \code{AttackerSkip}). Lemmas of the auxiliary
models are verified in \code{lightning_aux.spthy}, one model per preprocessor flag
(\code{-D=PTLC}, \code{-D=HTLC}, \code{-D=CLOCK}, \code{-D=CLTV}, \code{-D=TIMEOUT}). Every
listing is copied verbatim from these files, with the proof script omitted. The subsections
follow the groups of Table~\ref{tab:groups}: groups 1 to 9 first, then group 0 (channel layer
and keys). We use the same naming convention as the table: A for an all-traces safety property,
E for an exists-trace witness (reachability or non-vacuity).

All results were obtained with \tamarin{} 1.10 and are reproduced by \texttt{make all}.
\texttt{make verify} proves all 55 lemmas of the full model, i.e., with all nine restrictions in
place, including the liveness restrictions T1, T2 and T3 (32 all-traces safety lemmas, 13
hop-bypass lemmas, 9 exists-trace lemmas found automatically, and
\code{Wormhole_Fee_Theft_Reachable}, replayed from a saved 116-step proof). \texttt{make general}
proves the 24 lemmas of the any-length theorem in the model without the three liveness
restrictions. \texttt{make aux} proves the 21 lemmas of \code{lightning_aux.spthy}, and
\texttt{make necessity} removes one liveness restriction at a time (Table~\ref{tab:necessity}).

\emph{Running a single check.} A lemma is checked on its own by calling \tamarin{} on its file
with \texttt{-{}-prove=} followed by the lemma name, and with the flags of its mode; an auxiliary
model is proved as a whole with its flag and \texttt{-{}-prove}. For example, for a lemma of
\texttt{make verify}, a lemma of the general mode and the CLTV model:
{\footnotesize
\begin{verbatim}
make all
tamarin-prover lightning.spthy --prove=ForwardTimeGap
tamarin-prover lightning.spthy -D=GENERAL -D=NO_IMC \
  -D=NO_OTF -D=NO_RBT -D=NECESSITY \
  --prove=Any_Hop_Skip_Requires_Compromise
tamarin-prover lightning_aux.spthy -D=CLTV --prove
\end{verbatim}
}
\noindent A single check uses the preceding lemmas marked \code{[reuse]} as assumptions;
\texttt{make verify} and \texttt{make general} prove those lemmas too. \code{Wormhole_Fee_Theft_Reachable} is checked with \texttt{-{}-lemma=} in place of
\texttt{-{}-prove=}, which replays the proof saved in the file. The TIMEOUT model also needs
\texttt{-{}-heuristic=c} and \texttt{-{}-stop-on-trace=seqdfs}. If memory is short,
\texttt{make verify-one} checks the lemmas of \texttt{make verify} one per run, and \texttt{MEM}
sets the heap limit of \tamarin{} (6\,GB by default).

\emph{Action facts.} Throughout, action facts carry their intended meaning from
Sec.~\ref{sec:core}; $y = h(x)$ is a payment hash and \texttt{ptr} a channel.
{\raggedright
\begin{itemize}
  \item \texttt{InvoiceCreated(R, S, y)}: receiver \texttt{R} issued an invoice for $y$ to sender
    \texttt{S}. The invoice, the tuple of the tag \texttt{\textquotesingle invoice\textquotesingle},
    \texttt{R}, \texttt{S} and $y$, is signed by \texttt{R} and carries no amount.
  \item \texttt{SecretPreimage(R, x)} and \texttt{PreimageReleased(R, x)}: \texttt{R} generated
    the preimage $x$, and \texttt{R} released it.
  \item \texttt{HTLCOffered(S, F1, ptr, y)}: the sender \texttt{S} offered an HTLC on $y$ to its
    first forwarder \texttt{F1} on channel \texttt{ptr}. \texttt{OfferFor(S, R, ptr, y)} annotates
    the same step with the receiver \texttt{R} of the invoice.
  \item \texttt{HTLCAdded(ptr, y)}: an HTLC on $y$ was added to channel \texttt{ptr}, by an offer
    or by a forward.
  \item \texttt{HTLCForwarded(Me, Prev, Next, ptrIn, ptrOut, y)}: forwarder \texttt{Me}, holding
    an incoming HTLC from \texttt{Prev} on \texttt{ptrIn}, added an outgoing HTLC to \texttt{Next}
    on \texttt{ptrOut} (one rule for every hop). \texttt{FeeCharged(Me, y, \%fee)} is the fee
    \texttt{Me} charges at that forward.
  \item \texttt{HTLCFulfilled(R, Prev, ptr, y)}: the receiver redeemed its incoming HTLC.
  \item \texttt{IntermediaryClaimed(Me, Prev, ptrIn, x)} and
    \texttt{FeeEarned(Me, ptrIn, ptrOut, y, \%fee)}: forwarder \texttt{Me} redeemed its incoming
    HTLC with $x$, and earned its fee.
  \item \texttt{SenderSettled(S, F1, R, ptr, x)}: the sender settled its HTLC with $x$.
  \item \texttt{Redeemed(ptr, y)} and \texttt{Refunded(Up, ptr, y)}: the HTLC on \texttt{ptr} was
    redeemed; or it was refunded, after its timeout, to \texttt{Up}, the party that had offered
    it.
  \item \texttt{TimedOut(ptr)}: the HTLC timeout of channel \texttt{ptr} was reached.
  \item \texttt{SettledHop(ptrIn, ptrOut, y)} and \texttt{RefundHop(ptrIn, ptrOut, y)}:
    annotations; the hop from \texttt{ptrIn} to \texttt{ptrOut} settled, or had its outgoing HTLC
    refunded.
  \item \texttt{UsedFulfilled(ptr, x)}: annotation; a settlement consumed the fulfilment of the
    HTLC on \texttt{ptr}.
  \item \texttt{FulfilledProduced(ptr, x)}: a fulfilment for \texttt{ptr} was produced by an
    intermediary settlement or by a skip.
  \item \texttt{Skipped(ptr, x)}: annotation of \code{AttackerSkip}; compromised parties relayed
    $x$ onto \texttt{ptr} (wormhole).
  \item \texttt{ChannelOpen(A, B, ptr)}: \texttt{A} and \texttt{B} opened channel \texttt{ptr}.
  \item \texttt{KeyGen(A, ltk)} and \texttt{LtkCompromised(A)}: \texttt{A} generated its long-term
    key \texttt{ltk}; the long-term key of \texttt{A} was compromised.
  \item $K(x)$: the adversary knows the term $x$.
\end{itemize}
\par}
\noindent The action facts of the channel layer are explained in Appendix~\ref{app:channel}.

% ---------------------------------------------------------------------------------- group 1
\subsection{Invoice and Preimage (Group 1)}\label{app:invoice}

\paragraph{\texttt{Invoice\_Authenticates\_Settlement} (A)}
The sender can settle an HTLC only if the payment hash corresponds to an invoice the receiver
actually issued, at a strictly earlier time; the only way around this is for the receiver's
long-term key to have been compromised (strictly earlier as well). In short, a settlement is
legitimate only when a genuine invoice backs it; a settlement without such an invoice is
attributable to a compromised receiver.
\begin{spthy}
lemma Invoice_Authenticates_Settlement:
    "All S F1 R ptr x #i.
        SenderSettled(S, F1, R, ptr, x) @ #i
     ==> (Ex #j. InvoiceCreated(R, S, h(x)) @ #j & #j < #i)
       | (Ex #r. LtkCompromised(R) @ #r & #r < #i)"
\end{spthy}

\paragraph{\texttt{Forged\_Invoice\_Requires\_Key\_Compromise} (E)}
Reachability witness. There is a trace in which the sender settles an HTLC for receiver
\texttt{R} although \texttt{R} never created an invoice for $h(x)$; the long-term key of
\texttt{R} is compromised in this trace, and no forwarding takes place. Together with
\code{Invoice_Authenticates_Settlement}, which requires such a compromise to precede the
settlement, this shows that the compromise case of that lemma is not vacuous.
\begin{spthy}
lemma Forged_Invoice_Requires_Key_Compromise:
    exists-trace
    "Ex S F1 R ptr x #i #r.
        SenderSettled(S, F1, R, ptr, x) @ #i
      & LtkCompromised(R) @ #r
      & not (Ex #j. InvoiceCreated(R, S, h(x)) @ #j)
      & not (Ex Me Prev Next pi po y #u. HTLCForwarded(Me, Prev, Next, pi, po, y) @ #u)"
\end{spthy}

\paragraph{\texttt{Invoice\_Has\_Secret\_Preimage} (A)}
Every invoice is created together with the receiver's secret preimage of its hash.
\begin{spthy}
lemma Invoice_Has_Secret_Preimage [reuse]:
    "All R S x #i.
        InvoiceCreated(R, S, h(x)) @ #i
     ==> SecretPreimage(R, x) @ #i"
\end{spthy}

\paragraph{\texttt{Preimage\_Secret\_Until\_Released} (A)}
As long as the receiver holds a secret preimage and has not released it, the adversary does not
know it: the first moment the adversary can learn the preimage is strictly after the receiver
has released it. Formally, if the receiver generated a secret preimage and the adversary later
learns it, then the receiver must have released that preimage at a strictly earlier point in
the trace. In this model the preimage is never sent to the network: fulfilments travel between
hops as state facts, and \code{AttackerSkip} relays them without the adversary learning $x$.
Hence \code{Preimage_Never_Known} holds, and this lemma holds trivially.
\begin{spthy}
lemma Preimage_Secret_Until_Released:
    "All R x #i #j.
        SecretPreimage(R, x) @ #i
      & K(x) @ #j
      & #i < #j
     ==> Ex #k. PreimageReleased(R, x) @ #k & #k < #j"
\end{spthy}

\paragraph{\texttt{Preimage\_Never\_Known} (A)}
The adversary never learns a receiver's preimage ($x$ is never output in this model);
\code{Preimage_Secret_Until_Released} therefore holds trivially.
\begin{spthy}
lemma Preimage_Never_Known:
    "All R x #i. SecretPreimage(R, x) @ #i ==> not (Ex #j. K(x) @ #j)"
\end{spthy}

\paragraph{\texttt{Invoice\_Released\_Once} (A)}
The receiver releases a given preimage at most once: if the same receiver releases the same
preimage at two times in a trace, those two times coincide, i.e., there is a single release
event. This holds because \code{Receiver_Fulfill_HTLC} consumes the linear fact
\texttt{ReceiverInv(R, S, x)}, which \code{Receiver_Create_Invoice} creates once per fresh
preimage $x$.
\begin{spthy}
lemma Invoice_Released_Once [reuse]:
    "All R x #i #j.
        PreimageReleased(R, x) @ #i
      & PreimageReleased(R, x) @ #j
     ==> #i = #j"
\end{spthy}

% ---------------------------------------------------------------------------------- group 2
\subsection{Lock Causality (Group 2)}\label{app:lock}

\paragraph{\texttt{Forward\_Requires\_Incoming} (A)}
A forwarder adds an outgoing HTLC only if an HTLC on the same hash was added earlier to its
incoming channel, either offered to it by a sender or forwarded to it by its predecessor. The
lemma covers every hop of a route of any length and needs no honesty assumption. It is
Lemma~\ref{lem:causality}.
\begin{spthy}
lemma Forward_Requires_Incoming [use_induction]:
    "All Me Prev Next ptrIn ptrOut y #i.
        HTLCForwarded(Me, Prev, Next, ptrIn, ptrOut, y) @ #i
     ==> (Ex S #j. HTLCOffered(S, Me, ptrIn, y) @ #j & #j < #i)
       | (Ex P0 pin0 #j. HTLCForwarded(Prev, P0, Me, pin0, ptrIn, y) @ #j & #j < #i)"
\end{spthy}

\paragraph{\texttt{HTLC\_On\_Opened\_Channel} (A)}
An HTLC is only ever added on a channel that has been opened beforehand: whenever an HTLC is
added on a channel, two parties opened that channel at a strictly earlier time.
\code{HTLCAdded} is emitted by \code{Sender_Offer_HTLC} and \code{Forward_HTLC}.
\begin{spthy}
lemma HTLC_On_Opened_Channel [reuse, use_induction]:
    "All ptr y #i.
        HTLCAdded(ptr, y) @ #i
     ==> Ex X Y #j. ChannelOpen(X, Y, ptr) @ #j & #j < #i"
\end{spthy}

\paragraph{\texttt{Offer\_Requires\_Invoice} (A)}
An HTLC is offered for receiver \texttt{R} only against an invoice that \texttt{R} created
earlier, or the long-term key of \texttt{R} is compromised at some point of the trace.
\begin{spthy}
lemma Offer_Requires_Invoice:
    "All S R ptr y #o.
        OfferFor(S, R, ptr, y) @ #o
     ==> (Ex #j. InvoiceCreated(R, S, y) @ #j & #j < #o)
       | (Ex #r. LtkCompromised(R) @ #r)"
\end{spthy}

% ---------------------------------------------------------------------------------- group 3
\subsection{Release Causality (Group 3)}\label{app:release}

\paragraph{\texttt{Claim\_Requires\_Release} (A)}
An intermediary can claim its incoming HTLC with a preimage only after (strictly later than) the
receiver has released that preimage: the release of the preimage is a prerequisite for any
intermediary settlement.
\begin{spthy}
lemma Claim_Requires_Release [use_induction]:
    "All MidP InP ptrIn R x #i #j.
        IntermediaryClaimed(MidP, InP, ptrIn, x) @ #i
      & SecretPreimage(R, x) @ #j
     ==> Ex #k. PreimageReleased(R, x) @ #k & #k < #i"
\end{spthy}

\paragraph{\texttt{Settle\_Requires\_Receiver\_Release} (A)}
When the sender settles an HTLC on a genuine invoice, the receiver must have released the
corresponding preimage strictly earlier: a settlement never precedes the receiver's release of
the preimage.
\begin{spthy}
lemma Settle_Requires_Receiver_Release:
    "All S F1 R ptr x #i #j.
        SenderSettled(S, F1, R, ptr, x) @ #i
      & InvoiceCreated(R, S, h(x)) @ #j
     ==> Ex #k. PreimageReleased(R, x) @ #k & #k < #i"
\end{spthy}

\paragraph{\texttt{Redeem\_Requires\_Receiver\_Release} (A)}
Every redeemed HTLC on $y$ comes with a receiver's release of a preimage $x$ with $h(x) = y$.
\begin{spthy}
lemma Redeem_Requires_Receiver_Release [use_induction]:
    "All ptr y #i.
        Redeemed(ptr, y) @ #i
     ==> Ex R x #k. PreimageReleased(R, x) @ #k & h(x) = y"
\end{spthy}

\paragraph{\texttt{Fulfilled\_Requires\_Release} (A)}
Every fulfilment produced by an intermediary settlement or a skip is preceded by a receiver's
release of $x$.
\begin{spthy}
lemma Fulfilled_Requires_Release [reuse, use_induction, heuristic=S]:
    "All ptr x #i.
        FulfilledProduced(ptr, x) @ #i
     ==> Ex R #k. PreimageReleased(R, x) @ #k & #k < #i"
\end{spthy}

\paragraph{\texttt{Fulfilled\_Source} (A)}
A fulfilment consumed by a settlement stems from an earlier redemption of that HTLC or from an
earlier skip.
\begin{spthy}
lemma Fulfilled_Source [reuse]:
    "All p x #i. UsedFulfilled(p, x) @ #i
     ==> (Ex #j. Redeemed(p, h(x)) @ #j & #j < #i)
       | (Ex #j. Skipped(p, x) @ #j & #j < #i)"
\end{spthy}

% ---------------------------------------------------------------------------------- group 4
\subsection{Atomicity (Group 4)}\label{app:atomicity}

\paragraph{\texttt{Payment\_Atomicity\_Under\_Liveness} (A)}
If two forwarders \texttt{F1} and \texttt{F2} forward consecutive HTLCs on the same hash $y$
(the outgoing channel of \texttt{F1} is the incoming channel of \texttt{F2}), the HTLC on the
outgoing channel of \texttt{F2} is redeemed, and no party's long-term key is ever compromised,
then the HTLC on the incoming channel of \texttt{F1} is never refunded: in a run without any key
compromise, once the payment is redeemed downstream of two forwarders, the party upstream of
them cannot be refunded. The lemma is proved in the full model, i.e., under the liveness
restrictions; it also verifies with T1 removed (Table~\ref{tab:necessity}). It is
Lemma~\ref{lem:atomicity}.
\begin{spthy}
lemma Payment_Atomicity_Under_Liveness:
    "All F1 F2 Prev1 Next1 Prev2 Next2 ptrSF1 ptrF1F2 ptrF2R y #i #j #k.
        HTLCForwarded(F1, Prev1, Next1, ptrSF1, ptrF1F2, y) @ #i
      & HTLCForwarded(F2, Prev2, Next2, ptrF1F2, ptrF2R, y) @ #j
      & Redeemed(ptrF2R, y) @ #k
      & (not Ex P #l. LtkCompromised(P) @ #l)
     ==> not (Ex Up #m. Refunded(Up, ptrSF1, y) @ #m)"
\end{spthy}

\paragraph{\texttt{Settle\_Excludes\_Sender\_Refund} (A)}
If the sender has settled an HTLC with a preimage, then that same HTLC is never refunded:
settlement and refund of the sender's HTLC are mutually exclusive.
\begin{spthy}
lemma Settle_Excludes_Sender_Refund [heuristic={sender_pending}]:
    "All S F1 R ptr x #i.
        SenderSettled(S, F1, R, ptr, x) @ #i
     ==> not (Ex #j. Refunded(S, ptr, h(x)) @ #j)"
\end{spthy}

\paragraph{\texttt{OneOutcomePerHTLC} (A)}
No HTLC is both redeemed and refunded. Both outcomes emit the action \code{Old} for the same
channel and hash, which the restriction \code{OneResolutionPerHTLC} allows only once.
\begin{spthy}
lemma OneOutcomePerHTLC [reuse]:
    "All Up ptr y #i #j.
       Redeemed(ptr, y) @ #i & Refunded(Up, ptr, y) @ #j ==> F"
\end{spthy}

\paragraph{\texttt{Multihop\_Routing\_Reachable} (E)}
Non-vacuity witness. The routing phase of a payment over two forwarders is reachable: the sender
offers, both forwarders forward and the receiver fulfils, with no other offer or forward, at most
one HTLC per channel, and no intermediary claim or sender settlement.
\begin{spthy}
lemma Multihop_Routing_Reachable:
    exists-trace
    "Ex S F1 F2 R p0 p1 p2 x #t2 #t3 #t4 #t5.
        HTLCOffered(S, F1, p0, h(x)) @ #t2
      & HTLCForwarded(F1, S, F2, p0, p1, h(x)) @ #t3
      & HTLCForwarded(F2, F1, R, p1, p2, h(x)) @ #t4
      & HTLCFulfilled(R, F2, p2, h(x)) @ #t5
      & #t2 < #t3 & #t3 < #t4 & #t4 < #t5
      & not (Ex A B p y #u. HTLCOffered(A, B, p, y) @ #u & not (A = S & B = F1))
      & not (Ex Me Prev Next pi po y #u. HTLCForwarded(Me, Prev, Next, pi, po, y) @ #u
              & not (Me = F1 & Prev = S & Next = F2)
              & not (Me = F2 & Prev = F1 & Next = R))
      & not (Ex ptr y1 y2 #m #n. HTLCAdded(ptr, y1) @ #m & HTLCAdded(ptr, y2) @ #n & not (#m = #n))
      & not (Ex Me P q z #e. IntermediaryClaimed(Me, P, q, z) @ #e)
      & not (Ex A B C q z #e. SenderSettled(A, B, C, q, z) @ #e)"
\end{spthy}

\paragraph{\texttt{Sanity\_Honest\_Payment\_Completes} (E)}
Non-vacuity witness. An honest payment over one forwarder completes end to end: the offer, the
forward, the fulfilment, the forwarder's claim and the sender's settlement occur, with no key
compromise, at most one HTLC per channel, exactly this claim and this settlement, and no refund.
\begin{spthy}
lemma Sanity_Honest_Payment_Completes:
  exists-trace
  "Ex S F1 R p0 p1 x #o #a #r #c #s.
      HTLCOffered(S, F1, p0, h(x)) @ #o
    & HTLCForwarded(F1, S, R, p0, p1, h(x)) @ #a
    & HTLCFulfilled(R, F1, p1, h(x)) @ #r
    & IntermediaryClaimed(F1, S, p0, x) @ #c
    & SenderSettled(S, F1, R, p0, x) @ #s
    & not (Ex P #k. LtkCompromised(P) @ #k)
    & not (Ex ptr y1 y2 #m #n. HTLCAdded(ptr, y1) @ #m & HTLCAdded(ptr, y2) @ #n & not (#m = #n))
    & (All Me P q z #e. IntermediaryClaimed(Me, P, q, z) @ #e ==> #e = #c)
    & (All A B C q z #e. SenderSettled(A, B, C, q, z) @ #e ==> #e = #s)
    & not (Ex U q z #e. Refunded(U, q, z) @ #e)"
\end{spthy}

\paragraph{\texttt{Sanity\_Atomic\_Payment\_Completes} (E)}
Non-vacuity witness. An honest payment over two forwarders completes end to end: the offer, both
forwards, the fulfilment, both intermediary claims and the sender's settlement occur, with no
key compromise, at most one HTLC per channel, exactly these claims and this settlement, and no
refund. It shows that the premise of \code{Payment_Atomicity_Under_Liveness} is satisfiable.
\begin{spthy}
lemma Sanity_Atomic_Payment_Completes:
  exists-trace
  "Ex S F1 F2 R p0 p1 p2 x #o #a #b #r #c1 #c2 #s.
      HTLCOffered(S, F1, p0, h(x)) @ #o
    & HTLCForwarded(F1, S, F2, p0, p1, h(x)) @ #a
    & HTLCForwarded(F2, F1, R, p1, p2, h(x)) @ #b
    & HTLCFulfilled(R, F2, p2, h(x)) @ #r
    & IntermediaryClaimed(F2, F1, p1, x) @ #c2
    & IntermediaryClaimed(F1, S, p0, x) @ #c1
    & SenderSettled(S, F1, R, p0, x) @ #s
    & not (Ex P #k. LtkCompromised(P) @ #k)
    & not (Ex ptr y1 y2 #m #n. HTLCAdded(ptr, y1) @ #m & HTLCAdded(ptr, y2) @ #n & not (#m = #n))
    & (All Me P q z #e. IntermediaryClaimed(Me, P, q, z) @ #e ==> (#e = #c1 | #e = #c2))
    & (All A B C q z #e. SenderSettled(A, B, C, q, z) @ #e ==> #e = #s)
    & not (Ex U q z #e. Refunded(U, q, z) @ #e)"
\end{spthy}

\paragraph{\texttt{Distinct\_Parties\_Configuration} (E)}
Non-vacuity witness. Three channels among four distinct parties can be opened with no
compromise, update, publication, close, cheat or forward, and with at most one HTLC per
channel.
\begin{spthy}
lemma Distinct_Parties_Configuration:
    exists-trace
    "Ex S F1 F2 R ptrSF1 ptrF1F2 ptrF2R #t1 #t2 #t3.
        ChannelOpen(S, F1, ptrSF1) @ #t1
      & ChannelOpen(F1, F2, ptrF1F2) @ #t2
      & ChannelOpen(F2, R, ptrF2R) @ #t3
      & not (S = F1) & not (S = F2) & not (S = R)
      & not (F1 = F2) & not (F1 = R) & not (F2 = R)
      & not (Ex P #i. LtkCompromised(P) @ #i)
      & not (Ex P Q p n #j. StateUpdate(P, Q, p, n) @ #j)
      & not (Ex P n #k. PublishState(P, n) @ #k)
      & not (Ex P Q p n #l. Close(P, Q, p, n) @ #l)
      & not (Ex P n #m. Cheat(P, n) @ #m)
      & not (Ex Me Prev Next pi po y #n. HTLCForwarded(Me, Prev, Next, pi, po, y) @ #n)
      & not (Ex ptr y1 y2 #m #n. HTLCAdded(ptr, y1) @ #m & HTLCAdded(ptr, y2) @ #n & not (#m = #n))"
\end{spthy}

% ---------------------------------------------------------------------------------- group 5
\subsection{Fees (Group 5)}\label{app:fees}

Amounts are local to a hop. \code{Forward_HTLC} receives the incoming amount, the fee and the
outgoing amount as inputs and emits an \code{Eq} action stating that the incoming amount is the
outgoing amount plus the fee, which the \code{Equality} restriction enforces; on every hop, the
incoming amount therefore equals the outgoing amount plus the fee by construction
(Sec.~\ref{sec:core-amounts}). The offer, the invoice and the route carry no amount, so amounts
are not chained along the path and the model makes no end-to-end value statement. The fee
guarantee is the following lemma.

\paragraph{\texttt{Fee\_Only\_On\_Successful\_Forward} (A)}
A forwarder whose key is never compromised earns its fee only after its outgoing HTLC was
redeemed and after it forwarded on that hop. It is Lemma~\ref{lem:fee}.
\begin{spthy}
lemma Fee_Only_On_Successful_Forward [heuristic={route_info_first}]:
    "All Me ptrIn ptrOut y %fee #i.
        FeeEarned(Me, ptrIn, ptrOut, y, %fee) @ #i
      & not (Ex #c. LtkCompromised(Me) @ #c)
     ==> (Ex #j. Redeemed(ptrOut, y) @ #j & #j < #i)
       & (Ex Prev Next #f. HTLCForwarded(Me, Prev, Next, ptrIn, ptrOut, y) @ #f & #f < #i)"
\end{spthy}

% ---------------------------------------------------------------------------------- group 6
\subsection{Refund and Griefing (Group 6)}\label{app:refund}

\paragraph{\texttt{Refund\_Requires\_Timeout} (A)}
Every refund of an HTLC is preceded by the timeout of its channel: an HTLC is refunded only
after that channel has timed out.
\begin{spthy}
lemma Refund_Requires_Timeout [reuse]:
    "All Up ptr y #i.
        Refunded(Up, ptr, y) @ #i
     ==> Ex #j. TimedOut(ptr) @ #j & #j < #i"
\end{spthy}

\paragraph{\texttt{Loss\_Requires\_Inaction} (A)}
If a forwarder's outgoing HTLC is redeemed while its incoming HTLC is refunded, then its incoming
HTLC is never redeemed. An intermediary therefore loses funds only through its own inaction, by
failing to redeem its incoming HTLC before it times out, and never because the same HTLC was both
redeemed and refunded (the two outcomes are exclusive).
\begin{spthy}
lemma Loss_Requires_Inaction:
    "All P Prev Next Up ptrIn ptrOut y #i #j #k.
        HTLCForwarded(P, Prev, Next, ptrIn, ptrOut, y) @ #i
      & Redeemed(ptrOut, y) @ #j
      & Refunded(Up, ptrIn, y) @ #k
     ==> not (Ex #l. Redeemed(ptrIn, y) @ #l)"
\end{spthy}

\paragraph{\texttt{Refund\_Possible} (E)}
Reachability witness. A payment can be refunded: the receiver issues an invoice, the sender
offers an HTLC on its hash, and the sender's HTLC is refunded, with no settlement, no key
compromise, no forwarding, and at most one HTLC per channel. Since a refund requires a timeout
(\code{Refund_Requires_Timeout}), the sender's funds stay locked until then;
Sec.~\ref{sec:wormhole} reports this trace as reachable griefing of the sender.
\begin{spthy}
lemma Refund_Possible:
    exists-trace
    "Ex S F1 R p y #i #j #k.
        InvoiceCreated(R, S, y) @ #j
      & HTLCOffered(S, F1, p, y) @ #k
      & Refunded(S, p, y) @ #i
      & #j < #k & #k < #i
      & not (Ex A MidP B p2 x #l. SenderSettled(A, MidP, B, p2, x) @ #l)
      & not (Ex P #r. LtkCompromised(P) @ #r)
      & not (Ex Me Prev Next pi po y2 #t1. HTLCForwarded(Me, Prev, Next, pi, po, y2) @ #t1)
      & not (Ex ptr y1 y2 #m #n. HTLCAdded(ptr, y1) @ #m & HTLCAdded(ptr, y2) @ #n & not (#m = #n))"
\end{spthy}

% ---------------------------------------------------------------------------------- group 7
\subsection{Wormhole (Group 7)}\label{app:wormhole}

\paragraph{\texttt{Wormhole\_Fee\_Theft\_Reachable} (E)}
Reachability witness (attack). There is a trace on the route Alice $\rightarrow$ Att1
$\rightarrow$ Indiana $\rightarrow$ Att2 $\rightarrow$ Bob in which Bob issues the invoice and
redeems his incoming HTLC and Alice settles her HTLC, yet the outgoing HTLC of the intermediary
Indiana is refunded and Indiana earns no fee. Alice, Bob and Indiana are never compromised, Att1
and Att2 are distinct parties, and every channel carries at most one HTLC. This is the
wormhole~\cite{malavolta2019} of Fig.~\ref{fig:wormhole}: the preimage is relayed around Indiana
by rule \code{AttackerSkip}, which only two compromised parties on the same payment can apply.
\tamarin{} does not find this trace automatically: \texttt{make verify} replays a saved 116-step
proof. The statement, in the pretty-printed notation of \tamarin{}, is shown in
Fig.~\ref{fig:wormhole-lemma}. In the PTLC model of \code{lightning_aux.spthy}, the
corresponding wormhole around an honest hop is unreachable (\code{No_Wormhole_Around_Honest_Hop}),
while the HTLC baseline reaches it (\code{Wormhole_Around_Honest_Hop}); see
Appendix~\ref{app:aux}.

\begin{figure*}[!t]
\begin{spthy}[basicstyle=\ttfamily\footnotesize]
lemma Wormhole_Fee_Theft_Reachable:
  exists-trace
  "∃ Alice Bob Att1 Att2 Indiana ptrAliceAtt1 ptrAtt1Indiana
     ptrIndianaAtt2 ptrAtt2Bob x #t1 #t2 #t3 #t4 #t5 #t6 #t7 #t8.
    (((((((((((((((InvoiceCreated( Bob, Alice, h(x) ) @ #t1) ∧
                  (HTLCOffered( Alice, Att1, ptrAliceAtt1, h(x) ) @ #t2)) ∧
                 (HTLCForwarded( Att1, Alice, Indiana, ptrAliceAtt1, ptrAtt1Indiana,
                                 h(x)
                  ) @ #t3)) ∧
                (HTLCForwarded( Indiana, Att1, Att2, ptrAtt1Indiana,
                                ptrIndianaAtt2, h(x)
                 ) @ #t4)) ∧
               (HTLCForwarded( Att2, Indiana, Bob, ptrIndianaAtt2, ptrAtt2Bob,
                               h(x)
                ) @ #t5)) ∧
              (HTLCFulfilled( Bob, Att2, ptrAtt2Bob, h(x) ) @ #t6)) ∧
             (SenderSettled( Alice, Att1, Bob, ptrAliceAtt1, x ) @ #t7)) ∧
            (Refunded( Indiana, ptrIndianaAtt2, h(x) ) @ #t8)) ∧
           (¬(Indiana = Alice))) ∧
          (¬(Indiana = Bob))) ∧
         (¬(Att1 = Att2))) ∧
        (¬(∃ #c. LtkCompromised( Alice ) @ #c))) ∧
       (¬(∃ #c. LtkCompromised( Bob ) @ #c))) ∧
      (¬(∃ #c. LtkCompromised( Indiana ) @ #c))) ∧
     (¬(∃ ptr1 ptr2 yy %fee #k.
         FeeEarned( Indiana, ptr1, ptr2, yy, %fee ) @ #k))) ∧
    (¬(∃ ptr y1 y2 #m #n.
        ((HTLCAdded( ptr, y1 ) @ #m) ∧ (HTLCAdded( ptr, y2 ) @ #n)) ∧
        (¬(#m = #n))))"
\end{spthy}
  \caption{Statement of \texttt{Wormhole\_Fee\_Theft\_Reachable} in \texttt{lightning.spthy},
    in the pretty-printed notation of \tamarin{} (exists-trace; verified by replaying a saved
    116-step proof).}
  \label{fig:wormhole-lemma}
\end{figure*}

\paragraph{\texttt{Skipped\_Requires\_Compromise} (A)}
Every wormhole skip happens after some party's long-term key was compromised. Its statement is
Lemma~\ref{lem:skip} in Sec.~\ref{sec:wormhole}. It is also re-proved without the liveness
restrictions (\texttt{make general}).

% ---------------------------------------------------------------------------------- group 8
\subsection{Intermediary Safety and Timing (Group 8)}\label{app:safety}

\subsubsection{Intermediary Safety}

\paragraph{\texttt{Intermediary\_Never\_Loses\_Under\_Liveness} (A)}
If an intermediary's outgoing HTLC is redeemed while its incoming HTLC is refunded, then its
long-term key was compromised. Equivalently, an honest intermediary whose outgoing HTLC is
redeemed never has its incoming HTLC refunded. The stronger lemma below places the compromise
before the incoming timeout. This lemma is Lemma~\ref{lem:intermediary}; it also verifies with
T1 or T2 removed (Table~\ref{tab:necessity}).
\begin{spthy}
lemma Intermediary_Never_Loses_Under_Liveness:
    "All P Prev Next Up ptrIn ptrOut y #i #j #k.
        HTLCForwarded(P, Prev, Next, ptrIn, ptrOut, y) @ #i
      & Redeemed(ptrOut, y) @ #j
      & Refunded(Up, ptrIn, y) @ #k
     ==> Ex #l. LtkCompromised(P) @ #l"
\end{spthy}

\paragraph{\texttt{Intermediary\_Never\_Loses\_Strong} (A)}
A forwarder whose outgoing HTLC is redeemed while its incoming HTLC is refunded was compromised
before its incoming HTLC timed out. It is Lemma~\ref{lem:intermediary-strong}; it also verifies
with T1 or T2 removed (Table~\ref{tab:necessity}).
\begin{spthy}
lemma Intermediary_Never_Loses_Strong:
    "All P Prev Next Up ptrIn ptrOut y #i #j #k.
        HTLCForwarded(P, Prev, Next, ptrIn, ptrOut, y) @ #i
      & Redeemed(ptrOut, y) @ #j
      & Refunded(Up, ptrIn, y) @ #k
     ==> Ex #l #t. LtkCompromised(P) @ #l & TimedOut(ptrIn) @ #t & #l < #t"
\end{spthy}

\paragraph{\texttt{Sanity\_Compromised\_Intermediary\_Can\_Lose} (E)}
Non-vacuity witness for \code{Intermediary_Never_Loses_Under_Liveness}: a compromised forwarder
can lose (its outgoing HTLC is fulfilled while its incoming HTLC is refunded), with at most one
HTLC per channel, no intermediary claim and no sender settlement.
\begin{spthy}
lemma Sanity_Compromised_Intermediary_Can_Lose:
  exists-trace
  "Ex S F1 R p0 p1 x #o #a #r #u #k.
      HTLCOffered(S, F1, p0, h(x)) @ #o
    & HTLCForwarded(F1, S, R, p0, p1, h(x)) @ #a
    & HTLCFulfilled(R, F1, p1, h(x)) @ #r
    & Refunded(S, p0, h(x)) @ #u
    & LtkCompromised(F1) @ #k
    & not (Ex ptr y1 y2 #m #n. HTLCAdded(ptr, y1) @ #m & HTLCAdded(ptr, y2) @ #n & not (#m = #n))
    & not (Ex Me P q z #e. IntermediaryClaimed(Me, P, q, z) @ #e)
    & not (Ex A B C q z #e. SenderSettled(A, B, C, q, z) @ #e)"
\end{spthy}

\paragraph{\texttt{Necessity\_IMC\_Honest\_Intermediary\_Loses} (E)}
Checked with \code{-D=NO_IMC} and \code{-D=NECESSITY} (\texttt{make necessity}, 28 steps).
Without T3 an honest forwarder can lose: there is a trace in which its outgoing HTLC is
fulfilled by the receiver while its incoming HTLC is refunded to the sender, with no key
compromise, no claim and no settlement.
Hence T3 is necessary for intermediary safety (Sec.~\ref{sec:timing}).
\begin{spthy}
lemma Necessity_IMC_Honest_Intermediary_Loses:
  exists-trace
  "Ex S F1 R p0 p1 x #o #a #r #u.
      HTLCOffered(S, F1, p0, h(x)) @ #o
    & HTLCForwarded(F1, S, R, p0, p1, h(x)) @ #a
    & HTLCFulfilled(R, F1, p1, h(x)) @ #r
    & Refunded(S, p0, h(x)) @ #u
    & not (Ex P #k. LtkCompromised(P) @ #k)
    & not (Ex ptr y1 y2 #m #n. HTLCAdded(ptr, y1) @ #m & HTLCAdded(ptr, y2) @ #n & not (#m = #n))
    & not (Ex Me P q z #e. IntermediaryClaimed(Me, P, q, z) @ #e)
    & not (Ex A B C q z #e. SenderSettled(A, B, C, q, z) @ #e)"
\end{spthy}

\subsubsection{Timing}

\paragraph{\texttt{Timeout\_Race\_Blocked} (A)}
No intermediary that is not compromised before its incoming timeout has its incoming HTLC
refunded when that timeout happens strictly before its outgoing HTLC is redeemed. The lemma also
verifies with any one of T1, T2 and T3 removed (4 steps each, Table~\ref{tab:necessity}), so it
does not depend on any single liveness restriction. In the abstract TIMEOUT model of
\code{lightning_aux.spthy}, where either channel may time out at any point and the outgoing HTLC
may be redeemed at any time, the race is reachable (\code{Early_Timeout_Race}, E, 6 steps;
Appendix~\ref{app:aux}).
\begin{spthy}
lemma Timeout_Race_Blocked:
    "not (Ex P Prev Next Up ptrIn ptrOut y #l #i #j #k.
        HTLCForwarded(P, Prev, Next, ptrIn, ptrOut, y) @ #l
      & TimedOut(ptrIn)     @ #k
      & Refunded(Up, ptrIn, y)  @ #j
      & Redeemed(ptrOut, y) @ #i
      & #k < #i
      & not (Ex #c. LtkCompromised(P) @ #c & #c < #k))"
\end{spthy}

\paragraph{\texttt{ForwardTimeGap} (A)}
If both channels of a forward time out and the forwarder is not compromised before its incoming
timeout, the outgoing HTLC times out strictly before the incoming one. The order of a hop's two
timeouts is thus not assumed: it is this lemma, derived from T2 (Sec.~\ref{sec:restrictions}).
It also verifies with T1 or T3 removed (Table~\ref{tab:necessity}). With \code{-D=NECESSITY},
the same statement is declared without \code{[reuse]}, so that no other proof uses it in the
necessity runs and in the general mode.
\begin{spthy}
lemma ForwardTimeGap [reuse]:
    "All P Prev Next ptrIn ptrOut y #f #j #k.
        HTLCForwarded(P, Prev, Next, ptrIn, ptrOut, y) @ #f
      & TimedOut(ptrOut) @ #j
      & TimedOut(ptrIn)  @ #k
      & not (Ex #c. LtkCompromised(P) @ #c & #c < #k)
     ==> #j < #k"
\end{spthy}

% ---------------------------------------------------------------------------------- group 9
\subsection{Hop Bypass (Group 9)}\label{app:bypass}

If the payment settles at the sender while a forwarder's own hop is refunded (its outgoing HTLC
is refunded, \code{RefundHop}), some party's key was compromised. The conclusion does not say
which party, or when. The lemmas for positions 1 to 6 are proved by \texttt{make verify}, in the
full model under T1 to T3, and need no single-offer premise; \code{Any_Hop_Skip_Requires_Compromise} covers any position on a route
of any length under a single-offer premise (every offer on $h(x)$ uses the sender's channel
\texttt{pSF}) and is proved by \texttt{make general}. Sec.~\ref{sec:bypass} discusses these
results.

\subsubsection{Positions 1 to 6}

\paragraph{\texttt{First\_Hop\_Skip\_Requires\_Compromise} (A)}
The first forwarder cannot be bypassed without a compromise: if the sender settles its HTLC on
channel \texttt{pSF} while the hop of its first forwarder \texttt{F1} from \texttt{pSF} is
refunded, some party's key was compromised.
\begin{spthy}
lemma First_Hop_Skip_Requires_Compromise [heuristic={lemmas_first}]:
    "All S F1 R pSF Next pOut x #s #f #r.
        SenderSettled(S, F1, R, pSF, x) @ #s
      & HTLCForwarded(F1, S, Next, pSF, pOut, h(x)) @ #f
      & RefundHop(pSF, pOut, h(x)) @ #r
     ==> Ex P #c. LtkCompromised(P) @ #c"
\end{spthy}

\paragraph{\texttt{Second\_Hop\_Skip\_Requires\_Compromise} (A)}
The second forwarder (the position of Indiana in \code{Wormhole_Fee_Theft_Reachable}) cannot be
bypassed without a compromise.
\begin{spthy}
lemma Second_Hop_Skip_Requires_Compromise [heuristic={lemmas_first}]:
    "All S F1 R pSF x F2 p1 Next pOut #s #f1 #f2 #r.
        SenderSettled(S, F1, R, pSF, x) @ #s
      & HTLCForwarded(F1, S, F2, pSF, p1, h(x)) @ #f1
      & HTLCForwarded(F2, F1, Next, p1, pOut, h(x)) @ #f2
      & RefundHop(p1, pOut, h(x)) @ #r
     ==> Ex P #c. LtkCompromised(P) @ #c"
\end{spthy}

\paragraph{\texttt{Hop3\_Skip\_Requires\_Compromise} (A)}
The forwarder at position 3 cannot be bypassed without a compromise. The statement adds a third
forward to the chain of \code{Second_Hop_Skip_Requires_Compromise}; the statements for positions
4 to 6 each add one more.
\begin{spthy}
lemma Hop3_Skip_Requires_Compromise [heuristic={lemmas_first}]:
    "All S R x pSF pOut Next F1 F2 F3 p1 p2 #s #r #f1 #f2 #f3.
        SenderSettled(S, F1, R, pSF, x) @ #s
      & HTLCForwarded(F1, S, F2, pSF, p1, h(x)) @ #f1
      & HTLCForwarded(F2, F1, F3, p1, p2, h(x)) @ #f2
      & HTLCForwarded(F3, F2, Next, p2, pOut, h(x)) @ #f3
      & RefundHop(p2, pOut, h(x)) @ #r
     ==> Ex P #c. LtkCompromised(P) @ #c"
\end{spthy}

\paragraph{\texttt{Hop4\_Skip\_Requires\_Compromise} (A)}
The forwarder at position 4 cannot be bypassed without a compromise.
\begin{spthy}
lemma Hop4_Skip_Requires_Compromise [heuristic={lemmas_first}]:
    "All S R x pSF pOut Next F1 F2 F3 F4 p1 p2 p3 #s #r #f1 #f2 #f3 #f4.
        SenderSettled(S, F1, R, pSF, x) @ #s
      & HTLCForwarded(F1, S, F2, pSF, p1, h(x)) @ #f1
      & HTLCForwarded(F2, F1, F3, p1, p2, h(x)) @ #f2
      & HTLCForwarded(F3, F2, F4, p2, p3, h(x)) @ #f3
      & HTLCForwarded(F4, F3, Next, p3, pOut, h(x)) @ #f4
      & RefundHop(p3, pOut, h(x)) @ #r
     ==> Ex P #c. LtkCompromised(P) @ #c"
\end{spthy}

\paragraph{\texttt{Hop5\_Skip\_Requires\_Compromise} (A)}
The forwarder at position 5 cannot be bypassed without a compromise.
\begin{spthy}
lemma Hop5_Skip_Requires_Compromise [heuristic={lemmas_first}]:
    "All S R x pSF pOut Next F1 F2 F3 F4 F5 p1 p2 p3 p4 #s #r #f1 #f2 #f3 #f4 #f5.
        SenderSettled(S, F1, R, pSF, x) @ #s
      & HTLCForwarded(F1, S, F2, pSF, p1, h(x)) @ #f1
      & HTLCForwarded(F2, F1, F3, p1, p2, h(x)) @ #f2
      & HTLCForwarded(F3, F2, F4, p2, p3, h(x)) @ #f3
      & HTLCForwarded(F4, F3, F5, p3, p4, h(x)) @ #f4
      & HTLCForwarded(F5, F4, Next, p4, pOut, h(x)) @ #f5
      & RefundHop(p4, pOut, h(x)) @ #r
     ==> Ex P #c. LtkCompromised(P) @ #c"
\end{spthy}

\paragraph{\texttt{Hop6\_Skip\_Requires\_Compromise} (A)}
The forwarder at position 6 cannot be bypassed without a compromise.
\begin{spthy}
lemma Hop6_Skip_Requires_Compromise [heuristic={lemmas_first}]:
    "All S R x pSF pOut Next F1 F2 F3 F4 F5 F6 p1 p2 p3 p4 p5 #s #r #f1 #f2 #f3 #f4 #f5 #f6.
        SenderSettled(S, F1, R, pSF, x) @ #s
      & HTLCForwarded(F1, S, F2, pSF, p1, h(x)) @ #f1
      & HTLCForwarded(F2, F1, F3, p1, p2, h(x)) @ #f2
      & HTLCForwarded(F3, F2, F4, p2, p3, h(x)) @ #f3
      & HTLCForwarded(F4, F3, F5, p3, p4, h(x)) @ #f4
      & HTLCForwarded(F5, F4, F6, p4, p5, h(x)) @ #f5
      & HTLCForwarded(F6, F5, Next, p5, pOut, h(x)) @ #f6
      & RefundHop(p5, pOut, h(x)) @ #r
     ==> Ex P #c. LtkCompromised(P) @ #c"
\end{spthy}

\subsubsection{Helper Lemmas}
The hop-bypass lemmas rest on seven helper lemmas, all of kind A and declared \code{[reuse]}.
All seven are also re-proved in \texttt{make general}.

\paragraph{\texttt{Forward\_Unique\_Per\_Incoming} (A)}
An incoming HTLC is forwarded at most once: two forwards of the same hash from the same incoming
channel are the same event.
\begin{spthy}
lemma Forward_Unique_Per_Incoming [reuse, heuristic={route_first}]:
    "All M1 P1 N1 M2 P2 N2 pIn o1 o2 y #i #j.
        HTLCForwarded(M1, P1, N1, pIn, o1, y) @ #i
      & HTLCForwarded(M2, P2, N2, pIn, o2, y) @ #j
     ==> #i = #j"
\end{spthy}

\paragraph{\texttt{SettledHop\_From\_Forward} (A)}
A hop settles only if it was forwarded earlier, from the same incoming to the same outgoing
channel.
\begin{spthy}
lemma SettledHop_From_Forward [reuse, heuristic={route_first}]:
    "All pIn pOut y #i. SettledHop(pIn, pOut, y) @ #i
     ==> Ex M P N #j. HTLCForwarded(M, P, N, pIn, pOut, y) @ #j & #j < #i"
\end{spthy}

\paragraph{\texttt{No\_Settle\_And\_Refund\_Same\_Hop} (A)}
No hop both settles and has its outgoing HTLC refunded.
\begin{spthy}
lemma No_Settle_And_Refund_Same_Hop [reuse, heuristic={route_first}]:
    "All pIn pOut y #i #j. SettledHop(pIn, pOut, y) @ #i & RefundHop(pIn, pOut, y) @ #j ==> F"
\end{spthy}

\paragraph{\texttt{No\_Fulfill\_On\_Forwarded\_Incoming} (A)}
The HTLC on a channel is never both fulfilled by a receiver and forwarded onwards.
\begin{spthy}
lemma No_Fulfill_On_Forwarded_Incoming [reuse, heuristic={route_first}]:
    "All R Q M P N p o y #i #j.
        HTLCFulfilled(R, Q, p, y) @ #i & HTLCForwarded(M, P, N, p, o, y) @ #j ==> F"
\end{spthy}

\paragraph{\texttt{Redeemed\_Source} (A)}
An HTLC is redeemed only by the settlement of the hop for which it is the incoming HTLC, or by
the receiver's fulfilment, in the same step.
\begin{spthy}
lemma Redeemed_Source [reuse]:
    "All p y #i. Redeemed(p, y) @ #i
     ==> (Ex o. SettledHop(p, o, y) @ #i)
       | (Ex R Q. HTLCFulfilled(R, Q, p, y) @ #i)"
\end{spthy}

\paragraph{\texttt{SenderSettled\_Uses\_Fulfilled} (A)}
A sender settlement consumes the fulfilment of the sender's own channel.
\begin{spthy}
lemma SenderSettled_Uses_Fulfilled [reuse]:
    "All S F1 R p x #i. SenderSettled(S, F1, R, p, x) @ #i ==> UsedFulfilled(p, x) @ #i"
\end{spthy}

\paragraph{\texttt{SettledHop\_Uses\_Outgoing} (A)}
A hop settlement consumes the fulfilment of the hop's outgoing channel, with a preimage of the
hop's hash.
\begin{spthy}
lemma SettledHop_Uses_Outgoing [reuse]:
    "All pIn pOut y #i. SettledHop(pIn, pOut, y) @ #i
     ==> Ex x. UsedFulfilled(pOut, x) @ #i & h(x) = y"
\end{spthy}

\subsubsection{Routes of Any Length (General Mode)}

\paragraph{\texttt{Any\_Hop\_Skip\_Requires\_Compromise} (A)}
Proved by \texttt{make general} in 7 steps, with the flags shown at the start of this
appendix. A single-offer payment that settles while any
forwarder's own hop is refunded requires a compromise. The lemma is proved in the model without
the three liveness restrictions; removing restrictions only adds traces, so the statement also
holds in the full model. It is Theorem~\ref{thm:anyhop}; the statement is written in the
pretty-printed notation of \tamarin{}.
\begin{spthy}
lemma Any_Hop_Skip_Requires_Compromise [heuristic={lemmas_first}]:
  all-traces
  "∀ S F1 R pSF x Me Prev Next pIn pOut #s #f #r.
    ((((SenderSettled( S, F1, R, pSF, x ) @ #s) ∧
       (HTLCForwarded( Me, Prev, Next, pIn, pOut, h(x) ) @ #f)) ∧
      (RefundHop( pIn, pOut, h(x) ) @ #r)) ∧
     (∀ A B p #o. (HTLCOffered( A, B, p, h(x) ) @ #o) ⇒ (p = pSF))) ⇒
    (∃ P #c. LtkCompromised( P ) @ #c)"
\end{spthy}

\noindent The general theorem uses a chain of seven helper lemmas, all of kind A and proved only
in this mode; we list them by name.
{\raggedright
\begin{itemize}
  \item \code{Forward_Parent}: the statement of \code{Forward_Requires_Incoming}, declared
    \code{[reuse]} for this mode.
  \item \code{Hop_Redeem_Propagates}: with no compromise, if the incoming HTLC of a forwarder is
    redeemed, so is its outgoing HTLC.
  \item \code{Sender_Channel_Redeemed}: with no compromise, a sender settlement implies that the
    sender's channel was redeemed.
  \item \code{Root_First}: every forward either forwards an offered HTLC or is preceded by a
    forward that does.
  \item \code{Next_Is_Child}: with a single offer, every forward on the hash that follows a
    forward to \texttt{Me} is, or comes after, the forward of \texttt{Me} from that channel.
  \item \code{Route_Redeemed_Gen}: with a single offer and no compromise, if the incoming HTLC of
    a forward is redeemed, so is the incoming HTLC of every forward on the same hash at the same
    time or later.
  \item \code{Route_Redeemed_No_Compromise}: with a single offer and no compromise, if the
    sender settles, the incoming HTLC of every forwarder on the hash is redeemed.
\end{itemize}
\par}
{\raggedright
\texttt{make general} proves 24 lemmas: the theorem, these seven helpers, and 16 lemmas of
\texttt{make verify} re-proved in the same mode, namely the seven helper lemmas of the hop-bypass
family above, \code{OneOutcomePerHTLC}, \code{Ltk_Known_Implies_Compromised},
\code{Invoice_Released_Once}, \code{Invoice_Has_Secret_Preimage}, \code{HTLC_On_Opened_Channel},
\code{Fulfilled_Requires_Release}, \code{Refund_Requires_Timeout}, \code{Fulfilled_Source} and
\code{Skipped_Requires_Compromise}.\par}

% ---------------------------------------------------------------------------------- group 0
\subsection{Channel Layer and Keys (Group 0)}\label{app:channel}

These lemmas concern the channel layer (Sec.~\ref{sec:core-channels}) and the long-term keys.
Their action facts read as follows.
{\raggedright
\begin{itemize}
  \item \texttt{StateUpdate(A, B, ptr, n)}: the channel moved to state \texttt{n} after both
    parties revoked their previous state.
  \item \texttt{NewStateBuilt(A, B, ptr, n, fA, fB)}: state \texttt{n} with balances
    \texttt{fA} and \texttt{fB} was built, at opening or by \code{Update_State}.
  \item \texttt{FundsLocked(A, B, ptr)}: the funds of both parties were locked when the channel
    opened.
  \item \texttt{PublishState(A, n)}: \texttt{A} published its commitment of state \texttt{n}.
  \item \texttt{Cheat(A, n)} and \texttt{Punish(B, A, n)}: \texttt{A} published its revoked
    state \texttt{n}, and \texttt{B} punished this and took all funds of the channel.
  \item \texttt{Close(A, B, ptr, n)}: the channel was closed cooperatively in state \texttt{n}.
  \item \texttt{OnChainDelayedProduce(A, n)} and \texttt{OnChainInstantProduce(A, n)}: an
    on-chain transaction of state \texttt{n} created a delayed (CSV-locked) or an instant output
    for \texttt{A}.
  \item \texttt{ClaimFunds(A, n)} and \texttt{FundsSettled(A, n)}: \texttt{A} claimed its delayed
    output after the CSV delay, or received its instant output.
  \item \texttt{SettledTx(A, tx)}: the transaction \texttt{tx} paid \texttt{A}.
\end{itemize}
\par}

\paragraph{\texttt{state\_update} (A)}
A channel is updated only after the same two parties opened it.
\begin{spthy}
lemma state_update [use_induction]:
    "All A B ptr n #i.
       StateUpdate(A, B, ptr, n) @ #i
     ==> Ex #j. ChannelOpen(A, B, ptr) @ #j & #j < #i"
\end{spthy}

\paragraph{\texttt{Funds\_Locked\_Before\_Update} (A)}
A channel is updated only after the funds of both parties were locked in it.
\begin{spthy}
lemma Funds_Locked_Before_Update [use_induction]:
    "All A B ptr n #i.
       StateUpdate(A, B, ptr, n) @ #i
     ==> Ex #j. FundsLocked(A, B, ptr) @ #j & #j < #i"
\end{spthy}

\paragraph{\texttt{Update\_Requires\_Negotiation} (A)}
A channel moves to a new state only after both parties built that state, with its balances.
\begin{spthy}
lemma Update_Requires_Negotiation [use_induction]:
    "All A B ptr n #i.
         StateUpdate(A, B, ptr, n) @ #i
     ==> Ex fA fB #j. NewStateBuilt(A, B, ptr, n, fA, fB) @ #j & #j < #i"
\end{spthy}

\paragraph{\texttt{delayed\_funds} (A)}
A party claims a delayed output of state \texttt{n} only after an on-chain transaction of state
\texttt{n} created that output for it.
\begin{spthy}
lemma delayed_funds:
    "All A n #i.
       ClaimFunds(A, n) @ #i
     ==> Ex #j. OnChainDelayedProduce(A, n) @ #j & #j < #i"
\end{spthy}

\paragraph{\texttt{instant\_funds} (A)}
A party receives an instant output of state \texttt{n} only after an on-chain transaction of
state \texttt{n} created that output for it.
\begin{spthy}
lemma instant_funds:
    "All A n #i.
       FundsSettled(A, n) @ #i
     ==> Ex #j. OnChainInstantProduce(A, n) @ #j & #j < #i"
\end{spthy}

\paragraph{\texttt{settlement\_is\_traceable} (A)}
Every instant settlement names the transaction that pays the party, in the same step.
\begin{spthy}
lemma settlement_is_traceable:
    "All A n #i.
       FundsSettled(A, n) @ #i
     ==> Ex tx #j. SettledTx(A, tx) @ #j & #j = #i"
\end{spthy}

\paragraph{\texttt{No\_Punishment\_Without\_Cheating} (A)}
A party is punished for state \texttt{n} only if it published that revoked state earlier, or if
its long-term key was compromised earlier.
\begin{spthy}
lemma No_Punishment_Without_Cheating:
    "All A B n #i.
        Punish(B, A, n) @ #i
     ==> (Ex #k. Cheat(A, n) @ #k & #k < #i)
       | (Ex #j. LtkCompromised(A) @ #j & #j < #i)"
\end{spthy}

\paragraph{\texttt{Ltk\_Known\_Implies\_Compromised} (A)}
The adversary knows a long-term key only if that key was compromised. Here \texttt{!KU(ltk)}
means that the adversary can construct \texttt{ltk}.
\begin{spthy}
lemma Ltk_Known_Implies_Compromised [reuse]:
    "All A ltk #i #j.
        KeyGen(A, ltk) @ #i
      & !KU(ltk) @ #j
     ==> Ex #k. LtkCompromised(A) @ #k"
\end{spthy}

\paragraph{\texttt{Protocol\_execution} (E)}
Non-vacuity witness. A channel can be opened, a new state built and committed by an update, and
the channel closed cooperatively in that state, in this order, with at most one HTLC per
channel.
\begin{spthy}
lemma Protocol_execution:
    exists-trace
    "Ex A B ptr n fA fB #i #j #k #l.
        ChannelOpen(A, B, ptr) @ #i
      & NewStateBuilt(A, B, ptr, n, fA, fB) @ #j
      & StateUpdate(A, B, ptr, n) @ #k
      & Close(A, B, ptr, n) @ #l
      & #i < #j & #j < #k & #k < #l
      & not (Ex ptr y1 y2 #m #n. HTLCAdded(ptr, y1) @ #m & HTLCAdded(ptr, y2) @ #n & not (#m = #n))"
\end{spthy}

\paragraph{\texttt{Cooperative\_Close\_Execution} (E)}
Non-vacuity witness. A channel can be opened and closed cooperatively without any state update,
with at most one HTLC per channel.
\begin{spthy}
lemma Cooperative_Close_Execution:
    exists-trace
    "Ex A B ptr n #k #i.
        ChannelOpen(A, B, ptr) @ #k & Close(A, B, ptr, n) @ #i & #k < #i
      & not (Ex A2 B2 ptr2 n2 #j. StateUpdate(A2, B2, ptr2, n2) @ #j)
      & not (Ex ptr y1 y2 #m #n. HTLCAdded(ptr, y1) @ #m & HTLCAdded(ptr, y2) @ #n & not (#m = #n))"
\end{spthy}

% ---------------------------------------------------------------------------------- aux
\subsection{Auxiliary Models}\label{app:aux}

The auxiliary models of \code{lightning_aux.spthy} are self-contained and are not composed with
\code{lightning.spthy} (Sec.~\ref{sec:structure}). Table~\ref{tab:aux} lists their 21 lemmas,
all verified by \texttt{make aux}; we print the key statements below.

The PTLC model (\code{-D=PTLC}) has the fixed route S $\rightarrow$ H1 $\rightarrow$ H2
$\rightarrow$ H3 $\rightarrow$ R over the channels \texttt{c0} to \texttt{c3}. Each channel has
its own secret, derived from the secret of the next channel by the hash chain
$s_{\mathit{in}} = h(\langle s_{\mathit{out}}, t\rangle)$ with a fresh tweak $t$ per hop, a
symbolic stand-in for the adaptor-signature points of~\cite{malavolta2019}; locks are hashes of
the secrets. The HTLC baseline (\code{-D=HTLC}) has the same route with one secret $x$ for every
channel. In both models a compromised hop leaks its tweak, the secret used to claim a channel
becomes public, and a compromised party can claim its incoming channel with any value the
adversary knows. There is no skip rule: a wormhole must be built from the adversary's knowledge.
The CLOCK model (\code{-D=CLOCK}) gives channels concrete CLTV deadlines on a block clock, lets a
forwarder forward only if its outgoing deadline is strictly lower than its incoming one, and
records the block heights at which HTLCs are redeemed, claimed and refunded. The CLTV model
(\code{-D=CLTV}) registers hops whose incoming deadline exceeds the outgoing one by a positive
delta (restriction \code{DeltaPositive}) and reveals a preimage only strictly before the
outgoing deadline (restriction \code{FulfillInTime}). The TIMEOUT model (\code{-D=TIMEOUT}) has four rules and no ordering:
either channel may time out at any point, and the outgoing HTLC may be redeemed at any time.

\begin{table}[!t]
  \caption{Lemmas of the auxiliary models (\texttt{lightning\_aux.spthy}, \texttt{make aux},
    21/21 verified). A: all-traces; E: exists-trace; Steps: size of the proof found by
    \tamarin{}.}
  \label{tab:aux}
  \centering
  \footnotesize
  \renewcommand{\arraystretch}{1.1}
  \setlength{\tabcolsep}{4pt}
  \begin{tabular}{@{}lL{5.45cm}cr@{}}
    \toprule
    Flag & Lemma & Kind & Steps\\
    \midrule
    \code{PTLC} & \code{Honest_Payment_Completes} & E & 10\\
     & \code{Colluders_Present_Payment_Completes} & E & 14\\
     & \code{No_Wormhole_Around_Honest_Hop} & A & 25\\
     & \code{Control_Refund_Of_H2_Outgoing} & E & 6\\
     & \code{Control_Wormhole_If_H2_Compromised} & E & 18\\
    \addlinespace
    \code{HTLC} & \code{Honest_Payment_Completes} & E & 10\\
     & \code{Colluders_Present_Payment_Completes} & E & 8\\
     & \code{Wormhole_Around_Honest_Hop} & E & 12\\
    \addlinespace
    \code{CLOCK} & \code{HTLC_Needs_Open_Channel} & A & 36\\
     & \code{No_HTLC_After_Close} & A & 367\\
     & \code{Claim_Window_Exists} & A & 79\\
     & \code{Transitive_Preimage_Before_Upstream_Deadline} & A & 187\\
     & \code{Loss_Implies_Skipped_Claim} & A & 861\\
     & \code{Outcome_Exclusive} & A & 2\\
     & \code{Redeem_Reachable} & E & 7\\
     & \code{Refund_Reachable} & E & 7\\
     & \code{Honest_Flow_Possible} & E & 14\\
    \addlinespace
    \code{CLTV} & \code{CLTV_Gap_Is_Positive} & A & 2\\
     & \code{Claim_Window_Nonempty} & A & 2\\
     & \code{Staggered_Path_Safe} & A & 2\\
    \addlinespace
    \code{TIMEOUT} & \code{Early_Timeout_Race} & E & 6\\
    \bottomrule
  \end{tabular}
\end{table}

\paragraph{\texttt{No\_Wormhole\_Around\_Honest\_Hop} (A, PTLC)}
In the PTLC model, the channel \texttt{c0} of the sender is never redeemed while the outgoing
channel \texttt{c2} of an uncompromised H2 is refunded, unless the sender is compromised. The
result depends on the tweak of H2 staying secret: once H2 is compromised, the wormhole is
reachable (\code{Control_Wormhole_If_H2_Compromised}), and the refund of the outgoing channel of
an uncompromised H2 alone is reachable too (\code{Control_Refund_Of_H2_Outgoing}).
\begin{spthy}
lemma No_Wormhole_Around_Honest_Hop:
  all-traces
  "not (
Ex S H1 H2 H3 R c0 c1 c2 c3 #p #a #b.
      Pay(S, H1, H2, H3, R, c0, c1, c2, c3) @ #p
    & Redeemed(c0) @ #a
    & Refunded(c2) @ #b
    & not (Ex #k. Compromised(H2) @ #k)
    & not (Ex #k. Compromised(S) @ #k))"
\end{spthy}

\paragraph{\texttt{Wormhole\_Around\_Honest\_Hop} (E, HTLC)}
With HTLCs, that wormhole around an uncompromised H2 is reachable, with an uncompromised sender.
\begin{spthy}
lemma Wormhole_Around_Honest_Hop:
  exists-trace
  "Ex S H1 H2 H3 R c0 c1 c2 c3 #p #a #b.
      Pay(S, H1, H2, H3, R, c0, c1, c2, c3) @ #p
    & Redeemed(c0) @ #a
    & Refunded(c2) @ #b
    & not (Ex #k. Compromised(H2) @ #k)
    & not (Ex #k. Compromised(S) @ #k)"
\end{spthy}

\paragraph{\texttt{Claim\_Window\_Exists} (A, CLOCK)}
If a forwarder's outgoing HTLC is redeemed at block height \texttt{bOut} and its incoming HTLC is
refunded at height \texttt{bRef}, then \texttt{bOut} is strictly lower than \texttt{bRef}: the
forwarder had a non-empty window in which to claim.
\begin{spthy}
lemma Claim_Window_Exists:
    "All F pIn pOut dIn dOut y bOut bRef #f #r #u.
        Forwarded(F, pIn, pOut, dIn, dOut, y) @ #f
      & RedeemedAt(pOut, y, bOut) @ #r
      & RefundedAt(pIn,  y, bRef) @ #u
     ==> bOut << bRef"
\end{spthy}

\paragraph{\texttt{Loss\_Implies\_Skipped\_Claim} (A, CLOCK)}
A forwarder whose outgoing HTLC is redeemed while its incoming HTLC is refunded never claimed its
incoming HTLC.
\begin{spthy}
lemma Loss_Implies_Skipped_Claim:
    "All F pIn pOut dIn dOut y #f #r #u.
        Forwarded(F, pIn, pOut, dIn, dOut, y) @ #f
      & Redeemed(pOut, y) @ #r
      & Refunded(pIn,  y) @ #u
     ==> not (Ex b #c. ClaimedAt(pIn, y, b) @ #c)"
\end{spthy}

\paragraph{\texttt{Claim\_Window\_Nonempty} (A, CLTV)}
A preimage revealed before the outgoing deadline of a hop is revealed before its incoming
deadline.
\begin{spthy}
lemma Claim_Window_Nonempty [reuse]:
    "All ptrOut ptrIn d_out d_in b #i #j.
       HopRegistered(ptrOut, ptrIn, d_out, d_in) @ #i
     & Revealed(b, ptrOut, d_out) @ #j
     ==> b << d_in"
\end{spthy}

\paragraph{\texttt{Staggered\_Path\_Safe} (A, CLTV)}
Over two staggered hops, a preimage revealed before the tightest deadline is revealed before both
upstream deadlines, so both intermediaries have a non-empty claim window.
\begin{spthy}
lemma Staggered_Path_Safe:
    "All ptr1 ptr2 ptr2b ptr3 d1 d2 d3 b #i #j #k.
       HopRegistered(ptr1, ptr2, d1, d2) @ #i
     & HopRegistered(ptr2b, ptr3, d2, d3) @ #j
     & Revealed(b, ptr1, d1) @ #k
     ==> b << d2 & b << d3"
\end{spthy}

\paragraph{\texttt{Early\_Timeout\_Race} (E, TIMEOUT)}
Without any ordering of timeouts, a forwarder's incoming HTLC times out and is refunded, and its
outgoing HTLC is redeemed after that timeout. \code{Timeout_Race_Blocked} excludes this pattern
in \code{lightning.spthy} for forwarders that are not compromised before their incoming timeout.
\begin{spthy}
lemma Early_Timeout_Race:
    exists-trace
    "Ex P ptrIn ptrOut y #f #k #j #i.
        ForwardHTLC(P, ptrIn, ptrOut, y) @ #f
      & TimedOut(ptrIn)     @ #k
      & Refunded(ptrIn, y)  @ #j
      & Redeemed(ptrOut, y) @ #i
      & #k < #i"
\end{spthy}
